\documentclass[pdf]{beamer}

\usepackage[orientation=poster,size=custom,width=91.44,height=60.96,scale=1.2,debug]{beamerposter}

\mode<presentation>{
\usetheme{EnglishtownSmaller}
}

%\usepackage[utf8]{inputenc}
%\usepackage[urw-garamond]{mathdesign}
%\usepackage[T1]{fontenc}
%\usepackage{ebgaramond}
\usepackage{booktabs}
\usepackage{fancyref}
\usepackage{tabularx}
\usepackage{graphicx}
\usepackage{siunitx}
\sisetup{separate-uncertainty=true}
\usepackage{listings}
\usepackage{hyperref} % covered with beamer?
\hypersetup{
colorlinks=true,
urlcolor=blue,
linkcolor=blue
}
\usepackage[export]{adjustbox}
\usepackage{svg}
\usepackage{amsmath,amsthm,amssymb}
%\boldmath
% Uncomment next line to restore figure numbers if you like that sort of thing
\setbeamertemplate{caption}[numbered]

% For bibliographies
%\usepackage[backend=biber,style=apa]{biblatex}
%\addbibresource{references/sunrise.bib}
\usepackage{cite} % this is an IEEE conference
%\usepackage[round,authoryear]{natbib}
\usepackage{cleveref}

% IF YOU WANT WE CAN TRY TO CHANGE THE RED BANNER TO A STARFIELD IN MIDNGHT BLUE AND PURPLE||| Yes; that would be much better

\title{Analyzing solar plasma density via radio waves}
\author{Nikitas Komnenos, Pranav Arunkumar, Andrew Feinstein, Joseph Girone, \& Max Rothman}
\institute{Science \& Engineering Magnet Program and Manalapan High School, Englishtown, NJ 07726}
\date{\today}

\begin{document}
\begin{frame}{}
\begin{columns}[T,totalwidth=\textwidth]
\begin{column}{12in}
\begin{minipage}[t][\textheight]{\linewidth}
\begin{abstract}
% Our high school has been provided with a dual-dipole antenna and RSP1A receiver as part  of University of Michigan's sunRISE Ground Radio Lab project. In addition to providing observations for that project, we are studying electron density as a function of distance from the sun through examination of radio signal correlated to solar phenomena.
%Electron density of is a vital property of solar plasma that can be used to infer the behavior of various solar phenomena including coronal mass ejections(CMEs) and solar wind propagation, both of which are vital for near-earth satellites and global communication. This project aims to create a mathematical model to calculate the electron density as a function of distance from the sun by taking measurements of radio activity corresponding to solar phenomena.
% vital property - sounds like life sign and vital overused 
% infer the behavior - not sure what you mean here 
% a mathematical model does not take measurements of radio activity
Electron density plays a key role in coronal mass ejection (CME) and solar wind propagation, which can have major impacts on near-earth satellites and global communication. We are studying ways to estimate the electron density as a function of distance from the sun, using mathematical models based on radio measurements correlated with other observations of solar phenomena.
% AVOID ``radio activity'' it confuses with radioactivity
\end{abstract}

\vfill
% WHAT IS YOUR QUESTION
\begin{block}{Can we estimate electron density of solar plasma?}
The characteristic radio frequency of Langmuir waves produced by the solar plasma is given by
\begin{equation}
f_p = 9000 \sqrt{n_e}, \label{eq:1}
\end{equation}
where $f_p$ is in \unit{\hertz} and $n_e$ is the electron density in \unit{\per\centi\meter\cubed} \cite{White2023Extreme}.
\end{block}

\vfill

\begin{center}
% Remove the distracting overlays
\includegraphics[width=0.33\columnwidth]{figures/antenna-drawing.png}
\includegraphics[width=0.33\columnwidth]{figures/setup-annotated.png}
\end{center}
\begin{block}{Observations with a 20 MHz dual-dipole antenna}
\begin{itemize}
\item \qtyproduct{7.3x6.1}{\meter} dual-dipole antenna \& RSP1A receiver provided by University of Michigan SunRISE GRL project \cite{AkhavanTafti2025SunRISE}. 
\item Daytime observations in Monmouth County, NJ, w phase delay cable on southern antenna element.
%\item  % WHAT SOFTWARE TO RECORD AND GET SPECTROGRAM
%\begin{center}
%photo and drawing here
%\end{center}
\end{itemize}
\end{block}

\vspace*{\fill}
\end{minipage}
\end{column}




\begin{column}{12in}
\begin{minipage}[t][\textheight]{\linewidth}
% GIVE A HEADLINE DESCRIPTION OF HOW THE METHOD WORKS
% e.g. Do spike events (is that a thing) correlate with other measurements of electron density 
% this is prime real estate on poster so put a useful message here rather than the word Methods
\begin{block}{Recording and analyzing spectrographs to create an empirical model}

\vspace{1.25in}

\begin{center}
% Is this image AI generated?
% Does it use imagery you did not create?
\includegraphics[width=0.9\columnwidth]{figures/diagram.png} % TODO: cite this diagram as AI
\end{center}

\vspace{1.25in}

\begin{itemize}
\item Record spectrograph snapshots every \qty{3}{\hour} for \qty{4}{w}. Signals are recorded with a software-defined radio receiver (RSP1A; SDRPlay; Princes Risborough, UK) connected to a Linux laptop (Acer Travelmate P633-M).
\item Detect significant spike events and calculate $n_e$ using \cref{eq:1}
\item Find independent measurements of $n_e$ from various databases (NASA OMNIWeb, etc.) and calculate the differences.
\item Use models that predict $n_e$ as a function of distance, such as Newkirk  \cite{Soni2026Closed} or Baumbach-Aellen \cite{Benkevitch2011Method}, to estimate change in plasma density from near sun to near earth.
        % For the Newkirk Model, cite this link: https://link.springer.com/article/10.1007/s10509-026-04639-4 (there is a cite button, but I don't know what format)
        % For the Baumbach-Aellen model, cite this paper: https://arxiv.org/pdf/1110.2516
\item The predicted and actual differences will be used to compute a correction factor that will be used in addition to the densities and distances to be fitted to an empirical model.
\end{itemize}
\end{block}
\end{minipage}
\end{column}



\begin{column}{12in}
\begin{minipage}[t][\textheight]{\linewidth}
\begin{block}{Expected results (study newly underway, Fall 2026)}
\begin{center}
% Remove the distracting overlays
\includegraphics[width=0.9\columnwidth]{figures/example_graph.png} % TODO: replace this image IMMEDIATELY
\end{center}
We predict:
\begin{itemize}
\item The fitted model will follow some sort of exponential decay model;
\item Our calculated $n_e$ should be marginally higher than what is recorded in NASA databases such as Parker Solar Probe (PSP) and OMNIWeb, since our detected signals originate closer to the sun.
\end{itemize}
\end{block}

%\hfill\includegraphics[height=2in]{figures/se-logo-square.png}
\begin{block}{} %Acknowledgements} % Make this small
We thank C~Bennett, J~Komitas, Dr P~Wolniansky, Dr~Evangelista, University of Michigan, the SunRISE Project, Manalapan High School Science \& Engineering. Contact authors: \href{mailto:428parunkumar@frhsd.com}{\texttt{428parunkumar@frhsd.com}} and \href{mailto:428nkomnenos@frhsd.com}{\texttt{428nkomnenos@frhsd.com}}.

\begin{minipage}[t][]{11in}
\tiny
\bibliographystyle{IEEEtran}
\bibliography{references/sunrise.bib}
\end{minipage}
\end{block}
\end{minipage}
\end{column}
\end{columns}
\end{frame}
\end{document}
